\chapter{Tang Regulated Verse}

\section{The Tang dynasty and poetry}

The Tang (618--907) is the golden age of \emph{shī}. The \emph{Complete Tang Poems}
(\zht{全唐詩}, compiled 1705) contains nearly 50,000 poems by over 2,200 poets. Reasons
often given:
\begin{itemize}
  \item \textbf{Examinations.} From the late 7th century the \emph{jinshi} \zht{進士}
        examination included composing poems on set topics in set forms, so every educated
        man trained in verse.
  \item \textbf{Social function.} Poems were written at farewells, banquets, on visiting
        temples, as letters, as gifts: verse was the medium of elite social life.
  \item \textbf{Formal maturity.} The rules of ``recent-style'' verse were codified around
        700, giving poets a demanding and expressive instrument.
\end{itemize}

\begin{important}[Periods of Tang poetry (conventional)]
\begin{center}
\begin{tabular}{lll}
\toprule
Period & Dates & Poets \\
\midrule
Early Tang \zht{初唐} & 618--712 & Wang Bo, Chen Zi'ang, Shen Quanqi, Song Zhiwen \\
High Tang \zht{盛唐} & 712--765 & Wang Wei, Li Bai, Du Fu, Meng Haoran, Wang Changling \\
Mid Tang \zht{中唐} & 766--835 & Bai Juyi, Yuan Zhen, Han Yu, Li He \\
Late Tang \zht{晚唐} & 836--907 & Li Shangyin, Du Mu, Wen Tingyun \\
\bottomrule
\end{tabular}
\end{center}
The An Lushan rebellion (755--763) is the great dividing line between the confident High
Tang and the more troubled Mid Tang.
\end{important}

\section{Old style and recent style}

\begin{definition}[\emph{Gutishi} and \emph{jintishi}]
\begin{itemize}
  \item \textbf{Old-style verse} \zht{古體詩} \emph{gǔtǐshī}: any length, rhyme may change,
        no tonal rules, parallelism optional. Includes \emph{yuefu}-style songs. (Li Bai's
        ``Bring in the Wine'' is old style.)
  \item \textbf{Recent-style (regulated) verse} \zht{近體詩} \emph{jìntǐshī}: fixed line
        count, fixed line length (5 or 7), one level-tone rhyme, prescribed tonal pattern,
        required parallelism in the regulated octave.
\end{itemize}
\end{definition}

The forms of recent-style verse:

\begin{center}
\begin{tabular}{llll}
\toprule
Form & Lines & Characters per line & Total \\
\midrule
\zht{五言絕句} pentasyllabic quatrain & 4 & 5 & 20 \\
\zht{七言絕句} heptasyllabic quatrain & 4 & 7 & 28 \\
\zht{五言律詩} pentasyllabic regulated octave & 8 & 5 & 40 \\
\zht{七言律詩} heptasyllabic regulated octave & 8 & 7 & 56 \\
\zht{排律} extended regulated verse & 10+ & 5 or 7 & --- \\
\bottomrule
\end{tabular}
\end{center}

\begin{remark}
Shorthand: \emph{lǜshī} \zht{律詩} = regulated octave (8 lines); \emph{juéjù}
\zht{絕句} = quatrain (``cut-off lines''). A quatrain need not be strictly regulated (Li
Bai's ``Quiet Night Thoughts'' and Wang Wei's ``Deer Park'' bend the tonal rules), but
the most admired quatrains usually are.
\end{remark}

\section{Level and oblique tones}

Middle Chinese (the language of the Tang) had four tone categories:

\begin{center}
\begin{tabular}{lllll}
\toprule
Category & Name & Prosodic class & Example & Modern Mandarin \\
\midrule
level & \zht{平} \emph{píng} & level \zht{平} & \zht{山} & shān (1st) \\
rising & \zht{上} \emph{shǎng} & oblique \zht{仄} & \zht{海} & hǎi (3rd) \\
departing & \zht{去} \emph{qù} & oblique \zht{仄} & \zht{見} & jiàn (4th) \\
entering & \zht{入} \emph{rù} & oblique \zht{仄} & \zht{白} & bái (2nd!) \\
\bottomrule
\end{tabular}
\end{center}

\begin{important}[Level versus oblique]
For prosody, tones are simply split into two classes: \textbf{level} (\zht{平}
\emph{píng}) and \textbf{oblique} (\zht{仄} \emph{zè}, ``deflected'' = rising, departing,
entering). As a rough rule, Mandarin 1st and 2nd tones derive from the Middle Chinese level
tone and 3rd and 4th tones are oblique, \emph{except} the \textbf{entering tone}: syllables
that ended in -p, -t, -k (still preserved in Cantonese) have been scattered across all four
Mandarin tones. \zht{白} \emph{bái}, \zht{國} \emph{guó}, \zht{日} \emph{rì},
\zht{一} \emph{yī} are all oblique, though \zht{白}, \zht{國} and \zht{一} look ``level'' in
Mandarin. When in doubt, check Cantonese: \zht{白} \emph{baak}, \zht{國} \emph{gwok}.
\end{important}

\section{The rules of the regulated octave}

\begin{enumerate}
  \item \textbf{Within a line}, tones alternate in pairs. The decisive positions are the
        2nd and 4th characters (5-character line) or 2nd, 4th and 6th (7-character
        line): these must differ (2 $\neq$ 4, 4 $\neq$ 6).
  \item \textbf{Within a couplet}, the two lines are tonally \emph{opposite} (\zht{對}
        \emph{duì}) in the decisive positions.
  \item \textbf{Between couplets}, the first line of a couplet \emph{adheres} (\zht{粘}
        \emph{nián}) to the last line of the previous couplet: same tone in position 2.
  \item \textbf{Rhyme}: one rhyme throughout, in the \emph{level} tone, at the end of every
        even line (and optionally line 1).
  \item \textbf{Parallelism}: the 2nd and 3rd couplets (lines 3--4 and 5--6) must be parallel
        couplets; the first and last usually are not.
  \item \textbf{Prohibitions}: no ``three levels at the end'' (\zht{三平尾}); avoid an
        ``isolated level tone'' (\zht{孤平}) in a line.
\end{enumerate}

The mnemonic for seven-character lines is \zht{一三五不論，二四六分明}: ``1, 3, 5 don't
matter; 2, 4, 6 must be clear'' (an oversimplification, but useful).

\subsection{The four basic five-character lines}

\begin{center}
\begin{tabular}{llll}
\toprule
Label & Pattern & Ends in & Used where \\
\midrule
A & \zht{仄仄平平仄} & oblique & odd lines (non-rhyming) \\
B & \zht{平平仄仄平} & level & even lines (rhyming) \\
a & \zht{平平平仄仄} & oblique & odd lines \\
b & \zht{仄仄仄平平} & level & even lines (or a rhyming line 1) \\
\bottomrule
\end{tabular}
\end{center}

An octave whose first line has oblique 2nd position (``oblique-start'', \zht{仄起}) runs
A B a b A B a b. Below, white cells are level, grey cells oblique; a dashed outline marks
the position where variation is freely allowed.

\begin{center}
\begin{tikzpicture}[x=0.7cm,y=0.7cm, every node/.style={font=\small}]
\node[font=\scriptsize, gray] at (1,1) {1};
\node[font=\scriptsize, gray] at (2,1) {2};
\node[font=\scriptsize, gray] at (3,1) {3};
\node[font=\scriptsize, gray] at (4,1) {4};
\node[font=\scriptsize, gray] at (5,1) {5};
\filldraw[fill=gray!45, draw=black, dashed] (0.55,-0.45) rectangle (1.45,0.45);
\node at (1,0) {\zht{仄}};
\filldraw[fill=gray!45, draw=black] (1.55,-0.45) rectangle (2.45,0.45);
\node at (2,0) {\zht{仄}};
\filldraw[fill=white, draw=black] (2.55,-0.45) rectangle (3.45,0.45);
\node at (3,0) {\zht{平}};
\filldraw[fill=white, draw=black] (3.55,-0.45) rectangle (4.45,0.45);
\node at (4,0) {\zht{平}};
\filldraw[fill=gray!45, draw=black] (4.55,-0.45) rectangle (5.45,0.45);
\node at (5,0) {\zht{仄}};
\node[font=\scriptsize, anchor=east] at (0.4,0) {line 1};
\filldraw[fill=white, draw=black, dashed] (0.55,-1.45) rectangle (1.45,-0.55);
\node at (1,-1) {\zht{平}};
\filldraw[fill=white, draw=black] (1.55,-1.45) rectangle (2.45,-0.55);
\node at (2,-1) {\zht{平}};
\filldraw[fill=gray!45, draw=black] (2.55,-1.45) rectangle (3.45,-0.55);
\node at (3,-1) {\zht{仄}};
\filldraw[fill=gray!45, draw=black] (3.55,-1.45) rectangle (4.45,-0.55);
\node at (4,-1) {\zht{仄}};
\filldraw[fill=white, draw=black] (4.55,-1.45) rectangle (5.45,-0.55);
\node at (5,-1) {\zht{平}};
\node[font=\scriptsize, blue!60!black, anchor=west] at (5.6,-1) {rhyme};
\node[font=\scriptsize, anchor=east] at (0.4,-1) {line 2};
\filldraw[fill=white, draw=black, dashed] (0.55,-2.45) rectangle (1.45,-1.55);
\node at (1,-2) {\zht{平}};
\filldraw[fill=white, draw=black] (1.55,-2.45) rectangle (2.45,-1.55);
\node at (2,-2) {\zht{平}};
\filldraw[fill=white, draw=black] (2.55,-2.45) rectangle (3.45,-1.55);
\node at (3,-2) {\zht{平}};
\filldraw[fill=gray!45, draw=black] (3.55,-2.45) rectangle (4.45,-1.55);
\node at (4,-2) {\zht{仄}};
\filldraw[fill=gray!45, draw=black] (4.55,-2.45) rectangle (5.45,-1.55);
\node at (5,-2) {\zht{仄}};
\node[font=\scriptsize, anchor=east] at (0.4,-2) {line 3};
\filldraw[fill=gray!45, draw=black, dashed] (0.55,-3.45) rectangle (1.45,-2.55);
\node at (1,-3) {\zht{仄}};
\filldraw[fill=gray!45, draw=black] (1.55,-3.45) rectangle (2.45,-2.55);
\node at (2,-3) {\zht{仄}};
\filldraw[fill=gray!45, draw=black] (2.55,-3.45) rectangle (3.45,-2.55);
\node at (3,-3) {\zht{仄}};
\filldraw[fill=white, draw=black] (3.55,-3.45) rectangle (4.45,-2.55);
\node at (4,-3) {\zht{平}};
\filldraw[fill=white, draw=black] (4.55,-3.45) rectangle (5.45,-2.55);
\node at (5,-3) {\zht{平}};
\node[font=\scriptsize, blue!60!black, anchor=west] at (5.6,-3) {rhyme};
\node[font=\scriptsize, anchor=east] at (0.4,-3) {line 4};
\filldraw[fill=gray!45, draw=black, dashed] (0.55,-4.45) rectangle (1.45,-3.55);
\node at (1,-4) {\zht{仄}};
\filldraw[fill=gray!45, draw=black] (1.55,-4.45) rectangle (2.45,-3.55);
\node at (2,-4) {\zht{仄}};
\filldraw[fill=white, draw=black] (2.55,-4.45) rectangle (3.45,-3.55);
\node at (3,-4) {\zht{平}};
\filldraw[fill=white, draw=black] (3.55,-4.45) rectangle (4.45,-3.55);
\node at (4,-4) {\zht{平}};
\filldraw[fill=gray!45, draw=black] (4.55,-4.45) rectangle (5.45,-3.55);
\node at (5,-4) {\zht{仄}};
\node[font=\scriptsize, anchor=east] at (0.4,-4) {line 5};
\filldraw[fill=white, draw=black, dashed] (0.55,-5.45) rectangle (1.45,-4.55);
\node at (1,-5) {\zht{平}};
\filldraw[fill=white, draw=black] (1.55,-5.45) rectangle (2.45,-4.55);
\node at (2,-5) {\zht{平}};
\filldraw[fill=gray!45, draw=black] (2.55,-5.45) rectangle (3.45,-4.55);
\node at (3,-5) {\zht{仄}};
\filldraw[fill=gray!45, draw=black] (3.55,-5.45) rectangle (4.45,-4.55);
\node at (4,-5) {\zht{仄}};
\filldraw[fill=white, draw=black] (4.55,-5.45) rectangle (5.45,-4.55);
\node at (5,-5) {\zht{平}};
\node[font=\scriptsize, blue!60!black, anchor=west] at (5.6,-5) {rhyme};
\node[font=\scriptsize, anchor=east] at (0.4,-5) {line 6};
\filldraw[fill=white, draw=black, dashed] (0.55,-6.45) rectangle (1.45,-5.55);
\node at (1,-6) {\zht{平}};
\filldraw[fill=white, draw=black] (1.55,-6.45) rectangle (2.45,-5.55);
\node at (2,-6) {\zht{平}};
\filldraw[fill=white, draw=black] (2.55,-6.45) rectangle (3.45,-5.55);
\node at (3,-6) {\zht{平}};
\filldraw[fill=gray!45, draw=black] (3.55,-6.45) rectangle (4.45,-5.55);
\node at (4,-6) {\zht{仄}};
\filldraw[fill=gray!45, draw=black] (4.55,-6.45) rectangle (5.45,-5.55);
\node at (5,-6) {\zht{仄}};
\node[font=\scriptsize, anchor=east] at (0.4,-6) {line 7};
\filldraw[fill=gray!45, draw=black, dashed] (0.55,-7.45) rectangle (1.45,-6.55);
\node at (1,-7) {\zht{仄}};
\filldraw[fill=gray!45, draw=black] (1.55,-7.45) rectangle (2.45,-6.55);
\node at (2,-7) {\zht{仄}};
\filldraw[fill=gray!45, draw=black] (2.55,-7.45) rectangle (3.45,-6.55);
\node at (3,-7) {\zht{仄}};
\filldraw[fill=white, draw=black] (3.55,-7.45) rectangle (4.45,-6.55);
\node at (4,-7) {\zht{平}};
\filldraw[fill=white, draw=black] (4.55,-7.45) rectangle (5.45,-6.55);
\node at (5,-7) {\zht{平}};
\node[font=\scriptsize, blue!60!black, anchor=west] at (5.6,-7) {rhyme};
\node[font=\scriptsize, anchor=east] at (0.4,-7) {line 8};
\filldraw[fill=white] (0.6,-8.25) rectangle (0.9,-7.95); \node[font=\scriptsize, anchor=west] at (0.95,-8.1) {level \zht{平}};
\filldraw[fill=gray!45] (3.1,-8.25) rectangle (3.4,-7.95); \node[font=\scriptsize, anchor=west] at (3.45,-8.1) {oblique \zht{仄}};
\end{tikzpicture}
\end{center}

Check the rules on the grid: in each row, columns 2 and 4 differ; in each couplet (rows
1--2, 3--4, \ldots) column 2 flips; between couplets (rows 2--3, 4--5, 6--7) column 2 stays
the same (adhesion). The level-tone rhyme lands on every even line.

\subsection{The four couplets}

\begin{center}
\begin{tikzpicture}[cp/.style={draw, thick, minimum width=3.2cm, minimum height=0.8cm,
  font=\small, align=center}]
  \node[cp, fill=gray!10] (c1) at (0,3) {lines 1--2};
  \node[cp, fill=orange!20] (c2) at (0,2) {lines 3--4};
  \node[cp, fill=orange!20] (c3) at (0,1) {lines 5--6};
  \node[cp, fill=gray!10] (c4) at (0,0) {lines 7--8};
  \node[anchor=west, font=\small] at (2,3) {\zht{首聯} head couplet: occasion, setting};
  \node[anchor=west, font=\small] at (2,2) {\zht{頷聯} chin couplet: \textbf{parallel}, scene};
  \node[anchor=west, font=\small] at (2,1) {\zht{頸聯} neck couplet: \textbf{parallel}, scene/feeling};
  \node[anchor=west, font=\small] at (2,0) {\zht{尾聯} tail couplet: turn, close};
  \draw[thick] (-1.7,0.6) -- (-1.8,0.6) -- (-1.8,2.4) -- (-1.7,2.4);
  \node[anchor=east, font=\scriptsize, align=right] at (-2,1.5) {parallel\\ required};
\end{tikzpicture}
\end{center}

\begin{aside}
The body metaphor (head, chin, neck, tail) is traditional. A common movement in the
octave: open with the situation, expand it in two descriptive parallel couplets, then turn
inward to the poet's feeling in the last couplet. Later critics formalised the quatrain's
four lines as \zht{起承轉合}: \emph{begin, continue, turn, close}.
\end{aside}

\section{A regulated quatrain: ``Climbing Stork Tower''}

We met Wang Zhihuan's quatrain in Chapter 1. Here it is against the A B a b pattern.
Each character is shown with its tone class; a cell outlined in red differs from the ideal
pattern.

\begin{center}
\begin{tikzpicture}[x=0.8cm,y=0.8cm, every node/.style={font=\small}]
\node[font=\scriptsize, gray] at (1,1) {1};
\node[font=\scriptsize, gray] at (2,1) {2};
\node[font=\scriptsize, gray] at (3,1) {3};
\node[font=\scriptsize, gray] at (4,1) {4};
\node[font=\scriptsize, gray] at (5,1) {5};
\filldraw[fill=gray!45, draw=black, dashed] (0.55,-0.45) rectangle (1.45,0.45);
\node at (1,0) {\zht{白}};
\filldraw[fill=gray!45, draw=black] (1.55,-0.45) rectangle (2.45,0.45);
\node at (2,0) {\zht{日}};
\filldraw[fill=white, draw=black] (2.55,-0.45) rectangle (3.45,0.45);
\node at (3,0) {\zht{依}};
\filldraw[fill=white, draw=black] (3.55,-0.45) rectangle (4.45,0.45);
\node at (4,0) {\zht{山}};
\filldraw[fill=gray!45, draw=black] (4.55,-0.45) rectangle (5.45,0.45);
\node at (5,0) {\zht{盡}};
\filldraw[fill=white, draw=black, dashed] (0.55,-1.55) rectangle (1.45,-0.65);
\node at (1,-1.1) {\zht{黃}};
\filldraw[fill=white, draw=black] (1.55,-1.55) rectangle (2.45,-0.65);
\node at (2,-1.1) {\zht{河}};
\filldraw[fill=gray!45, draw=black] (2.55,-1.55) rectangle (3.45,-0.65);
\node at (3,-1.1) {\zht{入}};
\filldraw[fill=gray!45, draw=black] (3.55,-1.55) rectangle (4.45,-0.65);
\node at (4,-1.1) {\zht{海}};
\filldraw[fill=white, draw=black] (4.55,-1.55) rectangle (5.45,-0.65);
\node at (5,-1.1) {\zht{流}};
\node[font=\scriptsize, blue!60!black, anchor=west] at (5.6,-1.1) {rhyme};
\filldraw[fill=gray!45, draw=red!80!black, line width=1.6pt] (0.55,-2.65) rectangle (1.45,-1.75);
\node at (1,-2.2) {\zht{欲}};
\filldraw[fill=white, draw=black] (1.55,-2.65) rectangle (2.45,-1.75);
\node at (2,-2.2) {\zht{窮}};
\filldraw[fill=white, draw=black] (2.55,-2.65) rectangle (3.45,-1.75);
\node at (3,-2.2) {\zht{千}};
\filldraw[fill=gray!45, draw=black] (3.55,-2.65) rectangle (4.45,-1.75);
\node at (4,-2.2) {\zht{里}};
\filldraw[fill=gray!45, draw=black] (4.55,-2.65) rectangle (5.45,-1.75);
\node at (5,-2.2) {\zht{目}};
\filldraw[fill=gray!45, draw=black, dashed] (0.55,-3.75) rectangle (1.45,-2.85);
\node at (1,-3.3) {\zht{更}};
\filldraw[fill=gray!45, draw=black] (1.55,-3.75) rectangle (2.45,-2.85);
\node at (2,-3.3) {\zht{上}};
\filldraw[fill=gray!45, draw=black] (2.55,-3.75) rectangle (3.45,-2.85);
\node at (3,-3.3) {\zht{一}};
\filldraw[fill=white, draw=black] (3.55,-3.75) rectangle (4.45,-2.85);
\node at (4,-3.3) {\zht{層}};
\filldraw[fill=white, draw=black] (4.55,-3.75) rectangle (5.45,-2.85);
\node at (5,-3.3) {\zht{樓}};
\node[font=\scriptsize, blue!60!black, anchor=west] at (5.6,-3.3) {rhyme};
\filldraw[fill=white] (0.6,-4.55) rectangle (0.9,-4.25); \node[font=\scriptsize, anchor=west] at (0.95,-4.4) {level \zht{平}};
\filldraw[fill=gray!45] (3.1,-4.55) rectangle (3.4,-4.25); \node[font=\scriptsize, anchor=west] at (3.45,-4.4) {oblique \zht{仄}};
\end{tikzpicture}
\end{center}

\begin{itemize}
  \item The only deviation is \zht{欲} (oblique for level) in position 1 of line 3, which is
        a free position: the poem is perfectly regulated.
  \item Rhyme: \zht{流} \emph{liú} and \zht{樓} \emph{lóu} (the \zht{尤} rhyme group), both level.
  \item Note the entering-tone characters \zht{白}, \zht{日}, \zht{入}, \zht{目},
        \zht{一}: two of them (\zht{白} \emph{bái}, \zht{一} \emph{yī}) are not 3rd or 4th tone in Mandarin, so a modern reader would mistake them for level tones.
  \item Both couplets are parallel, which is not required in a quatrain; the
        \zht{起承轉合} movement is clear: scene, scene, turn (``if you want \ldots''),
        close (``climb'').
\end{itemize}

\section{Why the rules matter to a reader in translation}

Tonal patterning cannot be translated. But knowing it explains several things we
\emph{can} see in translation:
\begin{enumerate}
  \item \textbf{Density and balance.} The tonal and syntactic symmetry of the middle
        couplets makes them feel like two halves of one thought; the reader is trained to
        look for the relation (contrast, complement, cause).
  \item \textbf{The weight of the last couplet.} Since the middle is fixed in parallel, the
        unparallel ending feels like a release, often a question or a personal statement.
  \item \textbf{Virtuosity.} Poets were admired for making the form sound effortless.
        Du Fu, the master of the regulated octave, said: \zht{語不驚人死不休} ``If my
        words don't startle people, I won't rest even in death.''
\end{enumerate}

\begin{important}[Exam checklist for any Tang poem]
Count lines and characters $\Rightarrow$ name the form (\emph{wujue}, \emph{qijue},
\emph{wulü}, \emph{qilü}, or old style). Find the rhyme words. Find the parallel
couplets. Identify the turn.
\end{important}

\section*{Summary}
\begin{itemize}
  \item Tang = golden age of \emph{shi}; examinations and social use of verse; four
        periods (Early, High, Mid, Late); An Lushan rebellion as the divide.
  \item Old-style vs recent-style (regulated) verse; \emph{lüshi} (8 lines) and \emph{jueju}
        (4 lines), 5 or 7 characters.
  \item Level (\zht{平}) vs oblique (\zht{仄}) tones; entering tone is oblique even when
        Mandarin hides it.
  \item Rules: alternation within a line, opposition within a couplet, adhesion between
        couplets, level-tone rhyme on even lines, parallel middle couplets.
\end{itemize}
